\chapter*{Carta de Apresentação}

À Coordenadoria de Cursos Quadrimestrais da Escola Politécnica da USP

Caro(s) Sr(s).:

\par Este relatório (Modelagem Termodinâmica de Soluções de Interesse para a
Indústria de Alimentos) foi preparado como meu Relatório de Estágio do Módulo
E-1 para o Departamento de Engenharia Química da Escola Politécnica da USP,
orientado pelo Prof\textordmasculine. Dr\textordmasculine\ Pedro Pessoa.

Se trata de uma análise da aplicabilidade de diversos modelos, propostos na
literatura para a solução do problema em questão, à dados obtidos de soluções
reais de carboidratos e aminoácidos, com um ou mais solutos.

Esse estágio foi realizado na modalidade de Iniciação Científica/Tecnológica,
compreendendo a porção final de um projeto com um ano de duração; esse documento
se refere a todas as atividades realizadas durante o período, mas é possível
discriminar as atividades realizadas durante o módulo de estágio observando
o cronograma adotado.

Atenciosamente,

\includegraphics[width=0.2\textwidth]{signature.png}

Pedro Henrique Callil Soares

\newpage
